\subsection{Diagonal operators in sequence spaces}

Let I be a non-empty set. Recall (EVT, I, p. 4 and p. 5) that the vector space of bounded families \(x = (x_i)_{i \in \mathrm{I}}\) of elements of \(\mathbf{C}\) , equipped with the norm \(\| x \| = \sup_i |x_i|\) , is a Banach space denoted by \(\ell_{\mathbf{C}}^{\infty}(\mathrm{I})\) .

Let \(p\) be an element of \([1, +\infty[\) . The space \(\ell_{\mathbf{C}}^{p}(\mathrm{I})\) of families \(x = (x_{i})_{i \in \mathrm{I}}\) of elements of \(\mathbf{C}\) such that the family \((|x_{i}|^{p})\) is summable is a vector space on which the map \(x \mapsto \| x \| = (\sum_{i} |x_{i}|^{p})^{1/p}\) is a norm which makes it a Banach space. If one equips the space I with the discrete topology and with the measure for which \(\mu(\{x\}) = 1\) for all \(x \in \mathrm{I}\) , this space is none other than the Banach space \(\mathrm{L}_{\mathbf{C}}^{p}(\mathrm{I}, \mu)\) . When \(p = 1\) or p = 2, this notation thus coincides with that of EVT, I, p. 4 and EVT, V, p. 18 (cf. INT, IV, p. 141, §4, n° 1, example).

For families \(x = (x_{i})_{i \in \mathrm{I}}\) and \(y = (y_{i})_{i \in \mathrm{I}}\) of complex numbers, one denotes by xy the family \((x_{i}y_{i})_{i \in \mathrm{I}}\) .

PROPOSITION 1. --- Let \(\lambda = (\lambda_i)_{i\in I}\) be a bounded family of elements of C. Let \(p\) be an element of \([1, +\infty]\) . Let E be the Banach space \(\ell_{\mathbf{C}}^p (\mathrm{I})\) .

\begin{enumerate}
\def\labelenumi{\alph{enumi})}
\item
  For every \(x \in E\) , we have \(\lambda x \in E\) and the map u: \(x \mapsto \lambda x\) is a continuous endomorphism of E, whose spectrum in the Banach algebra \(\mathcal{L}(E)\) is equal to the closure in C of the set of \(\lambda_{i}\) , and whose norm is equal to \(\sup_{i} |\lambda_{i}|\) ;
\item
  The endomorphism u is compact if and only if for every real number \(\varepsilon > 0\) , the set of elements \(i \in I\) such that \(|\lambda_{i}| \geqslant \varepsilon\) is finite.
\end{enumerate}

With the notation of the proposition, the endomorphism \(x \mapsto \lambda x\) is called a diagonal endomorphism of \(\ell_{\mathbf{C}}^{p}(\mathrm{I})\) .

Let us prove the proposition. Let \(C = \sup_{i} |\lambda_{i}|\) . Let \(x \in E\) . We have the inequalities

\[
\sum_ {i \in I} | \lambda_ {i} x _ {i} | ^ {p} \leqslant C ^ {p} \| x \| ^ {p}, \text {   if   } p \neq + \infty
\]

\[
\sup _ {i \in I} | \lambda_ {i} x _ {i} | \leqslant C \| x \|, \text {   if   } p = + \infty .
\]

This proves that \(\lambda x \in E\) . The map \(x \mapsto \lambda x\) is therefore an endomorphism of E, and these inequalities prove that it has norm \(\leqslant C\) .

For \(j \in \mathrm{I}\) , denote by \(e_j\) the element \((x_i)_{i \in \mathrm{I}}\) of \(\ell_{\mathbf{C}}^p(\mathrm{I})\) such that \(x_j = 1\) and \(x_i = 0\) if \(i \neq j\) . We then have \(u(e_j) = \lambda_j e_j\) for all \(j \in \mathrm{I}\) , which shows that \(\lambda_j\) belongs to the spectrum of \(u\) . Since the spectrum of \(u\) is closed, the closure in \(\mathbf{C}\) of the set of \(\lambda_i\) is contained in the spectrum of \(u\) . Since the spectrum of \(u\) is contained in the disk centered at 0 with radius \(\|u\|\) (I, p. 24, th. 1 and formula (3) of I, p. 21), we deduce the inequality \(C \leqslant \|u\|\) , whence \(\|u\| = C\) .

Conversely, if \(\lambda \in \mathbf{C}\) is not adherent to the set of \(\lambda_{i}\) , then the family \(((\lambda_{i} - \lambda)^{-1})_{i\in \mathrm{I}}\) is bounded. The preceding argument therefore shows that the linear map \(x\mapsto ((\lambda -\lambda_i)^{-1}x_i)_{i\in \mathrm{I}}\) is an endomorphism of E. It is the inverse of \(u - \lambda 1_{\mathrm{E}}\) , and hence \(\lambda\) does not belong to the spectrum of \(u\) . This proves \(a\) ).

Let us now prove \(b\) ). Suppose that the endomorphism \(u\) is compact. Let \(\varepsilon > 0\) . Denote by J the set of \(i \in I\) such that \(|\lambda_i| \geqslant \varepsilon\) .

The set of elements \(e_i\) for \(i \in \mathbf{J}\) is bounded in E. Its image under \(u\) , consisting of the elements \(\lambda_i e_i\) for \(i \in \mathbf{J}\) is therefore relatively compact in E (Remark 1 of III, p. 2). Since \(\| \lambda_i e_i - \lambda_j e_j \| \geqslant \varepsilon\) for every pair of elements \(i \neq j\) in J, this is possible only if the set J is finite.

Let us prove the converse assertion. Let J be a finite subset of I and denote \(z_{\mathrm{J}} = (z_{\mathrm{J},i})_{i \in \mathrm{I}}\) , where \(z_{J,i} = \lambda_i\) for \(i \in J\) and \(z_{J,i} = 0\) for \(i \in I - J\) ; the family \(z_J\) defines a finite-rank endomorphism \(u_J \colon x \mapsto z_J x\) of E. By the preceding, the norm of the linear map \(u - u_J\) is bounded above by the supremum in \(R_+\) of the family \((|\lambda_i|)_{i \in I-J}\) . The hypothesis implies that for every \(\varepsilon > 0\) , there exists a finite subset \(J \subset I\) such that this upper bound is \(\leqslant \varepsilon\) . It follows that the map u is adherent to \(\mathcal{L}^{\mathrm{f}}(\mathrm{E})\) , and hence that it is compact (III, p. 4, cor. of prop. 2).

Remarks. --- 1) The hypothesis of assertion b) is always valid if I is finite. When the set I is infinite, it means that the family \((\lambda_{i})_{i\in\mathrm{I}}\) tends to 0 along the filter of complements of finite subsets of I; in particular, the set of \(i\in I\) such that \(\lambda_{i}\neq0\) is then countable.

\begin{enumerate}
\def\labelenumi{\arabic{enumi})}
\setcounter{enumi}{1}
\tightlist
\item
  We shall see later (cf. IV, p. 149, Theorem 1) that when \(p = 2\) every compact normal endomorphism of the Hilbertian space \(\ell_{\mathbf{C}}^{2}(\mathrm{I})\) is of the form \(u = w \circ v \circ w^{-1}\) where \(v\) is a compact diagonal endomorphism and \(w\) is a unitary endomorphism (EVT, V, p. 40) of the Hilbert space \(\ell_{\mathbf{C}}^{2}(\mathrm{I}).*\)
\end{enumerate}

